\par\noindent\begin{minipage}{\linewidth}\subsection*{KS01 / Limit universal property}
Grouped source statements. The universe and smallness conventions are retained.
\par\smallskip\noindent{\small\textbf{F} \path{ViewSections.lean}, line 39.\par\noindent\path{KomorebiSpace.ViewSystem.evaluations_jointly_injective}\par{\sffamily\footnotesize Source SHA-256 prefix: 53bb5c493e5cada2.}}\par
\par\smallskip\noindent{\small\textbf{F} \path{ViewSections.lean}, line 99.\par\noindent\path{KomorebiSpace.ViewSystem.Cone.exists_unique_lift}\par{\sffamily\footnotesize Source SHA-256 prefix: 53bb5c493e5cada2.}}\par
\smallskip\noindent{\color{KomorebiMuted}\rule{\linewidth}{.3pt}}\par\end{minipage}\par\medskip
\par\noindent\begin{minipage}{\linewidth}\subsection*{KS02 / Functorial transport}
Direct coordinate-level transport laws.
\par\smallskip\noindent{\small\textbf{F} \path{ViewSections.lean}, line 73.\par\noindent\path{KomorebiSpace.ViewSystem.Morphism.onAmbient_id}\par{\sffamily\footnotesize Source SHA-256 prefix: 53bb5c493e5cada2.}}\par
\par\smallskip\noindent{\small\textbf{F} \path{ViewSections.lean}, line 75.\par\noindent\path{KomorebiSpace.ViewSystem.Morphism.onAmbient_comp}\par{\sffamily\footnotesize Source SHA-256 prefix: 53bb5c493e5cada2.}}\par
\smallskip\noindent{\color{KomorebiMuted}\rule{\linewidth}{.3pt}}\par\end{minipage}\par\medskip
\par\noindent\begin{minipage}{\linewidth}\subsection*{KS03 / Terminal evaluation}
An already-global terminal fiber is recovered, not proved inhabited.
\par\smallskip\noindent{\small\textbf{F} \path{ViewSections.lean}, line 142.\par\noindent\path{KomorebiSpace.ViewSystem.terminalEvaluationEquiv}\par{\sffamily\footnotesize Source SHA-256 prefix: 53bb5c493e5cada2.}}\par
\smallskip\noindent{\color{KomorebiMuted}\rule{\linewidth}{.3pt}}\par\end{minipage}\par\medskip
\par\noindent\begin{minipage}{\linewidth}\subsection*{KS04 / Strict section gluing}
Requires a transition cover, not just object coverage.
\par\smallskip\noindent{\small\textbf{F} \path{SectionGluing.lean}, line 98.\par\noindent\path{KomorebiSpace.ViewSystem.gluingEquiv}\par{\sffamily\footnotesize Source SHA-256 prefix: cf412d554b19ca22.}}\par
\par\smallskip\noindent{\small\textbf{F} \path{SectionGluing.lean}, line 85.\par\noindent\path{KomorebiSpace.ViewSystem.Matching.exists_unique_glue}\par{\sffamily\footnotesize Source SHA-256 prefix: cf412d554b19ca22.}}\par
\smallskip\noindent{\color{KomorebiMuted}\rule{\linewidth}{.3pt}}\par\end{minipage}\par\medskip
\par\noindent\begin{minipage}{\linewidth}\subsection*{KS05 / Object-only cover failure}
The Boolean counterexample is implemented.
\par\smallskip\noindent{\small\textbf{F} \path{ViewBoundaryExamples.lean}, line 36.\par\noindent\path{KomorebiSpace.ViewBoundaryExamples.singleton_matching_not_gluable}\par{\sffamily\footnotesize Source SHA-256 prefix: c094eb03ea6b155f.}}\par
\smallskip\noindent{\color{KomorebiMuted}\rule{\linewidth}{.3pt}}\par\end{minipage}\par\medskip
\par\noindent\begin{minipage}{\linewidth}\subsection*{KS06 / Correspondence fibers}
Ordinary set fibers, not homotopy fibers.
\par\smallskip\noindent{\small\textbf{F} \path{ViewCorrespondence.lean}, line 43.\par\noindent\path{KomorebiSpace.ViewSystem.bridgeEquiv}\par{\sffamily\footnotesize Source SHA-256 prefix: a89331eabd7662f8.}}\par
\par\smallskip\noindent{\small\textbf{F} \path{ViewCorrespondence.lean}, line 94.\par\noindent\path{KomorebiSpace.ViewSystem.bridge_exists_iff_matching}\par{\sffamily\footnotesize Source SHA-256 prefix: a89331eabd7662f8.}}\par
\smallskip\noindent{\color{KomorebiMuted}\rule{\linewidth}{.3pt}}\par\end{minipage}\par\medskip
\par\noindent\begin{minipage}{\linewidth}\subsection*{KS07 / Nonempty fibers without a section}
The one-object Boolean swap has no global fixed choice.
\par\smallskip\noindent{\small\textbf{F} \path{ViewBoundaryExamples.lean}, line 90.\par\noindent\path{KomorebiSpace.ViewBoundaryExamples.swap_no_ambient}\par{\sffamily\footnotesize Source SHA-256 prefix: c094eb03ea6b155f.}}\par
\par\smallskip\noindent{\small\textbf{F} \path{ViewBoundaryExamples.lean}, line 87.\par\noindent\path{KomorebiSpace.ViewBoundaryExamples.swap_transports_bijective}\par{\sffamily\footnotesize Source SHA-256 prefix: c094eb03ea6b155f.}}\par
\smallskip\noindent{\color{KomorebiMuted}\rule{\linewidth}{.3pt}}\par\end{minipage}\par\medskip
\par\noindent\begin{minipage}{\linewidth}\subsection*{KS08 / Additive view differential}
Uses the actual arrow-indexed differential.
\par\smallskip\noindent{\small\textbf{F} \path{SectionObstruction.lean}, line 58.\par\noindent\path{KomorebiSpace.AdditiveViewSystem.differential_comp}\par{\sffamily\footnotesize Source SHA-256 prefix: 9089dd26b6bf34a0.}}\par
\par\smallskip\noindent{\small\textbf{F} \path{SectionObstruction.lean}, line 64.\par\noindent\path{KomorebiSpace.AdditiveViewSystem.differential_zero_iff}\par{\sffamily\footnotesize Source SHA-256 prefix: 9089dd26b6bf34a0.}}\par
\smallskip\noindent{\color{KomorebiMuted}\rule{\linewidth}{.3pt}}\par\end{minipage}\par\medskip
\par\noindent\begin{minipage}{\linewidth}\subsection*{KS09 / Additive extension criterion}
The quotient is defined independently of solvability.
\par\smallskip\noindent{\small\textbf{F} \path{SectionObstruction.lean}, line 126.\par\noindent\path{KomorebiSpace.AdditiveViewSystem.extensionClass_zero_iff}\par{\sffamily\footnotesize Source SHA-256 prefix: 9089dd26b6bf34a0.}}\par
\par\smallskip\noindent{\small\textbf{F} \path{SectionObstruction.lean}, line 142.\par\noindent\path{KomorebiSpace.AdditiveViewSystem.extensionClass_independent}\par{\sffamily\footnotesize Source SHA-256 prefix: 9089dd26b6bf34a0.}}\par
\smallskip\noindent{\color{KomorebiMuted}\rule{\linewidth}{.3pt}}\par\end{minipage}\par\medskip
\par\noindent\begin{minipage}{\linewidth}\subsection*{KS10 / Extension ambiguity torsor}
Translation requires one chosen extension.
\par\smallskip\noindent{\small\textbf{F} \path{SectionExtensionFibers.lean}, line 35.\par\noindent\path{KomorebiSpace.AdditiveViewSystem.extensionParameters}\par{\sffamily\footnotesize Source SHA-256 prefix: e422b1f42cdb28a1.}}\par
\par\smallskip\noindent{\small\textbf{F} \path{SectionExtensionFibers.lean}, line 62.\par\noindent\path{KomorebiSpace.AdditiveViewSystem.unique_extension_iff}\par{\sffamily\footnotesize Source SHA-256 prefix: e422b1f42cdb28a1.}}\par
\smallskip\noindent{\color{KomorebiMuted}\rule{\linewidth}{.3pt}}\par\end{minipage}\par\medskip
\par\noindent\begin{minipage}{\linewidth}\subsection*{KS11 / Parity seed example}
The disconnected free-coordinate variant in the text is an additional written consequence, not a separate Lean declaration.
\par\smallskip\noindent{\small\textbf{F} \path{SectionExtensionExample.lean}, line 65.\par\noindent\path{KomorebiSpace.SectionExtensionExample.extensionClass_zero_iff_even}\par{\sffamily\footnotesize Source SHA-256 prefix: 12f3be43a2401a41.}}\par
\par\smallskip\noindent{\small\textbf{F} \path{SectionExtensionExample.lean}, line 79.\par\noindent\path{KomorebiSpace.SectionExtensionExample.odd_seed_obstructed}\par{\sffamily\footnotesize Source SHA-256 prefix: 12f3be43a2401a41.}}\par
\par\smallskip\noindent{\small\textbf{F} \path{SectionExtensionExample.lean}, line 85.\par\noindent\path{KomorebiSpace.SectionExtensionExample.nonzero_seed_extends}\par{\sffamily\footnotesize Source SHA-256 prefix: 12f3be43a2401a41.}}\par
\smallskip\noindent{\color{KomorebiMuted}\rule{\linewidth}{.3pt}}\par\end{minipage}\par\medskip
\par\noindent\begin{minipage}{\linewidth}\subsection*{KS12 / Exactness at partial sections}
No surjectivity onto the whole obstruction quotient is asserted.
\par\smallskip\noindent{\small\textbf{F} \path{SectionExactSequence.lean}, line 108.\par\noindent\path{KomorebiSpace.AdditiveViewSystem.connecting_kernel_eq_range}\par{\sffamily\footnotesize Source SHA-256 prefix: 242e72f1d4e431ee.}}\par
\par\smallskip\noindent{\small\textbf{F} \path{SectionExactSequence.lean}, line 132.\par\noindent\path{KomorebiSpace.AdditiveViewSystem.restriction_surjective_iff_connecting_zero}\par{\sffamily\footnotesize Source SHA-256 prefix: 242e72f1d4e431ee.}}\par
\smallskip\noindent{\color{KomorebiMuted}\rule{\linewidth}{.3pt}}\par\end{minipage}\par\medskip
\par\noindent\begin{minipage}{\linewidth}\subsection*{KS13 / Extension by stages}
The dependent sum retains the intermediate choice.
\par\smallskip\noindent{\small\textbf{F} \path{SectionExtensionStages.lean}, line 28.\par\noindent\path{KomorebiSpace.ViewSystem.extensionStages}\par{\sffamily\footnotesize Source SHA-256 prefix: 8196e1272a7300f7.}}\par
\par\smallskip\noindent{\small\textbf{F} \path{SectionExtensionStages.lean}, line 50.\par\noindent\path{KomorebiSpace.ViewSystem.extension_nonempty_iff_intermediate}\par{\sffamily\footnotesize Source SHA-256 prefix: 8196e1272a7300f7.}}\par
\smallskip\noindent{\color{KomorebiMuted}\rule{\linewidth}{.3pt}}\par\end{minipage}\par\medskip
\par\noindent\begin{minipage}{\linewidth}\subsection*{KS14 / Protected correction}
The protected observable is a declared additive map.
\par\smallskip\noindent{\small\textbf{F} \path{Correction.lean}, line 42.\par\noindent\path{KomorebiSpace.CorrectionProblem.flatten_iff}\par{\sffamily\footnotesize Source SHA-256 prefix: 089d9155e0a43f66.}}\par
\par\smallskip\noindent{\small\textbf{F} \path{Correction.lean}, line 56.\par\noindent\path{KomorebiSpace.CorrectionProblem.obstruction_invariant}\par{\sffamily\footnotesize Source SHA-256 prefix: 089d9155e0a43f66.}}\par
\smallskip\noindent{\color{KomorebiMuted}\rule{\linewidth}{.3pt}}\par\end{minipage}\par\medskip
\par\noindent\begin{minipage}{\linewidth}\subsection*{KS15 / Obstruction witness}
A witness annihilates all admissible correction defects.
\par\smallskip\noindent{\small\textbf{F} \path{Correction.lean}, line 64.\par\noindent\path{KomorebiSpace.CorrectionProblem.nonzero_certificate}\par{\sffamily\footnotesize Source SHA-256 prefix: 089d9155e0a43f66.}}\par
\smallskip\noindent{\color{KomorebiMuted}\rule{\linewidth}{.3pt}}\par\end{minipage}\par\medskip
\par\noindent\begin{minipage}{\linewidth}\subsection*{KS16 / Exact tower maps and realized points}
The maps use a fixed sequential schedule. The record instances also supply identity and composition.
\par\smallskip\noindent{\small\textbf{F} \path{Towers.lean}, line 54.\par\noindent\path{KomorebiSpace.Tower.PreHom.exact_comp}\par{\sffamily\footnotesize Source SHA-256 prefix: fdd2819fff6150d2.}}\par
\par\smallskip\noindent{\small\textbf{F} \path{Realization.lean}, line 66.\par\noindent\path{KomorebiSpace.Resolution.Map.onLimit_comp}\par{\sffamily\footnotesize Source SHA-256 prefix: 90e0a22f2f7eebbe.}}\par
\smallskip\noindent{\color{KomorebiMuted}\rule{\linewidth}{.3pt}}\par\end{minipage}\par\medskip
\par\noindent\begin{minipage}{\linewidth}\subsection*{KS17 / Tower differential composition}
No additional mixed term is present in this differential.
\par\smallskip\noindent{\small\textbf{F} \path{Towers.lean}, line 75.\par\noindent\path{KomorebiSpace.Tower.PreHom.defect_comp}\par{\sffamily\footnotesize Source SHA-256 prefix: fdd2819fff6150d2.}}\par
\par\smallskip\noindent{\small\textbf{F} \path{TowerCorrection.lean}, line 39.\par\noindent\path{KomorebiSpace.TowerCorrection.exact_iff_differential_zero}\par{\sffamily\footnotesize Source SHA-256 prefix: 6710b485edbbb676.}}\par
\smallskip\noindent{\color{KomorebiMuted}\rule{\linewidth}{.3pt}}\par\end{minipage}\par\medskip
\par\noindent\begin{minipage}{\linewidth}\subsection*{KS18 / Cofinal point equivalence}
The subsequent alternating example is an explanatory specialization, not a statement that residual norms are invariant.
\par\smallskip\noindent{\small\textbf{F} \path{Cofinal.lean}, line 94.\par\noindent\path{KomorebiSpace.Resolution.cofinalEquiv}\par{\sffamily\footnotesize Source SHA-256 prefix: 0306f5a0a7f05e14.}}\par
\smallskip\noindent{\color{KomorebiMuted}\rule{\linewidth}{.3pt}}\par\end{minipage}\par\medskip
\par\noindent\begin{minipage}{\linewidth}\subsection*{KS19 / Geometric operator repair}
The correction map and residual contraction are hypotheses.
\par\smallskip\noindent{\small\textbf{F} \path{RepairDynamics.lean}, line 70.\par\noindent\path{KomorebiSpace.RepairDynamics.exactification_from_correction_operator}\par{\sffamily\footnotesize Source SHA-256 prefix: 419b196d29b2f14d.}}\par
\par\smallskip\noindent{\small\textbf{F} \path{RepairProjection.lean}, line 128.\par\noindent\path{KomorebiSpace.RepairDynamics.Scheme.projection_bound}\par{\sffamily\footnotesize Source SHA-256 prefix: 4292822995be802a.}}\par
\par\smallskip\noindent{\small\textbf{F} \path{RepairProjection.lean}, line 132.\par\noindent\path{KomorebiSpace.RepairDynamics.Scheme.projection_idempotent}\par{\sffamily\footnotesize Source SHA-256 prefix: 4292822995be802a.}}\par
\smallskip\noindent{\color{KomorebiMuted}\rule{\linewidth}{.3pt}}\par\end{minipage}\par\medskip
\par\noindent\begin{minipage}{\linewidth}\subsection*{KS20 / Constructed split tower correction}
Target splittings and the strict propagation gap are retained.
\par\smallskip\noindent{\small\textbf{F} \path{SplitTowerRepair.lean}, line 154.\par\noindent\path{KomorebiSpace.BanachTower.Splitting.differential_correction}\par{\sffamily\footnotesize Source SHA-256 prefix: 2dca3052b0338605.}}\par
\par\smallskip\noindent{\small\textbf{F} \path{SplitTowerRepair.lean}, line 199.\par\noindent\path{KomorebiSpace.BanachTower.Splitting.repair_error}\par{\sffamily\footnotesize Source SHA-256 prefix: 2dca3052b0338605.}}\par
\smallskip\noindent{\color{KomorebiMuted}\rule{\linewidth}{.3pt}}\par\end{minipage}\par\medskip
\par\noindent\begin{minipage}{\linewidth}\subsection*{KS21 / Split semantic and residual coordinates}
A continuous linear equivalence, not an isometry or an absolute semantic coordinate system.
\par\smallskip\noindent{\small\textbf{F} \path{SplitTowerCoordinates.lean}, line 98.\par\noindent\path{KomorebiSpace.BanachTower.Splitting.coordinates}\par{\sffamily\footnotesize Source SHA-256 prefix: c028e9936850deb9.}}\par
\par\smallskip\noindent{\small\textbf{F} \path{SplitTowerCoordinates.lean}, line 68.\par\noindent\path{KomorebiSpace.BanachTower.Splitting.coordinate_inverse}\par{\sffamily\footnotesize Source SHA-256 prefix: c028e9936850deb9.}}\par
\par\smallskip\noindent{\small\textbf{F} \path{SplitTowerSemantics.lean}, line 141.\par\noindent\path{KomorebiSpace.BanachTower.Splitting.residual_observable_injective}\par{\sffamily\footnotesize Source SHA-256 prefix: ecb64d5ac3d5369d.}}\par
\smallskip\noindent{\color{KomorebiMuted}\rule{\linewidth}{.3pt}}\par\end{minipage}\par\medskip
\par\noindent\begin{minipage}{\linewidth}\subsection*{KS22 / Growing identity repair}
The source bonding maps have the factor one half.
\par\smallskip\noindent{\small\textbf{F} \path{GrowingTower.lean}, line 195.\par\noindent\path{KomorebiSpace.GrowingTower.identityRepair_formula}\par{\sffamily\footnotesize Source SHA-256 prefix: 988f6ae0da7405cd.}}\par
\smallskip\noindent{\color{KomorebiMuted}\rule{\linewidth}{.3pt}}\par\end{minipage}\par\medskip
\par\noindent\begin{minipage}{\linewidth}\subsection*{KS23 / A bounded residual without bounded primitive}
This is a counterexample at the threshold, not necessity of the gap in every tower.
\par\smallskip\noindent{\small\textbf{F} \path{IdentityTowerObstruction.lean}, line 55.\par\noindent\path{KomorebiSpace.IdentityTowerObstruction.no_bounded_primitive}\par{\sffamily\footnotesize Source SHA-256 prefix: 026f9504c3c35507.}}\par
\par\smallskip\noindent{\small\textbf{F} \path{IdentityTowerObstruction.lean}, line 76.\par\noindent\path{KomorebiSpace.IdentityTowerObstruction.repair_exists_but_right_inverse_does_not}\par{\sffamily\footnotesize Source SHA-256 prefix: 026f9504c3c35507.}}\par
\smallskip\noindent{\color{KomorebiMuted}\rule{\linewidth}{.3pt}}\par\end{minipage}\par\medskip
\par\noindent\begin{minipage}{\linewidth}\subsection*{KS24 / Dependence on target splitting}
The two first-level matrices in the exposition are derived from the displayed source recurrence.
\par\smallskip\noindent{\small\textbf{F} \path{SplittingDependence.lean}, line 101.\par\noindent\path{KomorebiSpace.SplittingDependence.repair_operators_differ}\par{\sffamily\footnotesize Source SHA-256 prefix: b5aa369adbd67f9f.}}\par
\par\smallskip\noindent{\small\textbf{F} \path{SplittingDependence.lean}, line 112.\par\noindent\path{KomorebiSpace.SplittingDependence.protected_profiles_differ}\par{\sffamily\footnotesize Source SHA-256 prefix: b5aa369adbd67f9f.}}\par
\smallskip\noindent{\color{KomorebiMuted}\rule{\linewidth}{.3pt}}\par\end{minipage}\par\medskip
\par\noindent\begin{minipage}{\linewidth}\subsection*{KS25 / Lossless dyadic residual split}
Maps are the actual typed dyadic coarsener and its supplied-section lift.
\par\smallskip\noindent{\small\textbf{F} \path{DyadicResidualSplit.lean}, line 39.\par\noindent\path{KomorebiSpace.DyadicResidual.coarsen_pair}\par{\sffamily\footnotesize Source SHA-256 prefix: c4d40d1c7797e5ad.}}\par
\par\smallskip\noindent{\small\textbf{F} \path{DyadicResidualSplit.lean}, line 113.\par\noindent\path{KomorebiSpace.DyadicResidual.coarsen_lift}\par{\sffamily\footnotesize Source SHA-256 prefix: c4d40d1c7797e5ad.}}\par
\par\smallskip\noindent{\small\textbf{F} \path{DyadicResidualSplit.lean}, line 153.\par\noindent\path{KomorebiSpace.DyadicResidual.reconstruction}\par{\sffamily\footnotesize Source SHA-256 prefix: c4d40d1c7797e5ad.}}\par
\smallskip\noindent{\color{KomorebiMuted}\rule{\linewidth}{.3pt}}\par\end{minipage}\par\medskip
\par\noindent\begin{minipage}{\linewidth}\subsection*{KS26 / Cascade telescope and persistent part}
Generic split cascade. No norm convergence is inferred.
\par\smallskip\noindent{\small\textbf{F} \path{CascadeProjections.lean}, line 79.\par\noindent\path{KomorebiSpace.CoarseningCascade.remainder_nested_left}\par{\sffamily\footnotesize Source SHA-256 prefix: b2865b41336df871.}}\par
\par\smallskip\noindent{\small\textbf{F} \path{CascadeProjections.lean}, line 84.\par\noindent\path{KomorebiSpace.CoarseningCascade.remainder_nested_right}\par{\sffamily\footnotesize Source SHA-256 prefix: b2865b41336df871.}}\par
\par\smallskip\noindent{\small\textbf{F} \path{CascadeProjections.lean}, line 136.\par\noindent\path{KomorebiSpace.CoarseningCascade.persistent_iff_details_zero}\par{\sffamily\footnotesize Source SHA-256 prefix: b2865b41336df871.}}\par
\par\smallskip\noindent{\small\textbf{F} \path{CoarseningCascade.lean}, line 73.\par\noindent\path{KomorebiSpace.CoarseningCascade.reconstruction}\par{\sffamily\footnotesize Source SHA-256 prefix: 2719d5a3accb4dde.}}\par
\smallskip\noindent{\color{KomorebiMuted}\rule{\linewidth}{.3pt}}\par\end{minipage}\par\medskip
\par\noindent\begin{minipage}{\linewidth}\subsection*{KS27 / Recovery by density}
Uniform boundedness of the restored remainders is a hypothesis.
\par\smallskip\noindent{\small\textbf{F} \path{CascadeDensity.lean}, line 75.\par\noindent\path{KomorebiSpace.CoarseningCascade.recoverable_iff_closure_finiteCore}\par{\sffamily\footnotesize Source SHA-256 prefix: 11d90a5995354d99.}}\par
\par\smallskip\noindent{\small\textbf{F} \path{CascadeDensity.lean}, line 81.\par\noindent\path{KomorebiSpace.CoarseningCascade.all_recoverable_iff_dense}\par{\sffamily\footnotesize Source SHA-256 prefix: 11d90a5995354d99.}}\par
\smallskip\noindent{\color{KomorebiMuted}\rule{\linewidth}{.3pt}}\par\end{minipage}\par\medskip
\par\noindent\begin{minipage}{\linewidth}\subsection*{KS28 / Dyadic finite-coordinate recovery}
The printed proof is explicitly an outline. The source contains the dependent-index induction.
\par\smallskip\noindent{\small\textbf{F} \path{DyadicInfiniteReconstruction.lean}, line 71.\par\noindent\path{KomorebiSpace.DyadicResidual.System.infinite_details_injective}\par{\sffamily\footnotesize Source SHA-256 prefix: 0c81242c3c6c0715.}}\par
\par\smallskip\noindent{\small\textbf{F} \path{DyadicInfiniteReconstruction.lean}, line 201.\par\noindent\path{KomorebiSpace.DyadicResidual.System.reconstruction_coordinate}\par{\sffamily\footnotesize Source SHA-256 prefix: 0c81242c3c6c0715.}}\par
\smallskip\noindent{\color{KomorebiMuted}\rule{\linewidth}{.3pt}}\par\end{minipage}\par\medskip
\par\noindent\begin{minipage}{\linewidth}\subsection*{KS29 / Shift cascade and c0}
The scalar expanding-lift boundary is separately implemented in CascadeBoundaryExamples.
\par\smallskip\noindent{\small\textbf{F} \path{SequenceCascade.lean}, line 140.\par\noindent\path{KomorebiSpace.SequenceCascade.recoverable_iff_vanishes}\par{\sffamily\footnotesize Source SHA-256 prefix: 0a384fa460f07336.}}\par
\par\smallskip\noindent{\small\textbf{F} \path{SequenceCascade.lean}, line 194.\par\noindent\path{KomorebiSpace.SequenceCascade.ones_remainder_norm}\par{\sffamily\footnotesize Source SHA-256 prefix: 0a384fa460f07336.}}\par
\par\smallskip\noindent{\small\textbf{F} \path{SequenceCascade.lean}, line 270.\par\noindent\path{KomorebiSpace.SequenceCascade.geometric_reconstruction_error}\par{\sffamily\footnotesize Source SHA-256 prefix: 0a384fa460f07336.}}\par
\smallskip\noindent{\color{KomorebiMuted}\rule{\linewidth}{.3pt}}\par\end{minipage}\par\medskip
\par\noindent\begin{minipage}{\linewidth}\subsection*{KS30 / Gauge groupoid and section transport}
The action and groupoid instances are also read in this module. Composition order is retained.
\par\smallskip\noindent{\small\textbf{F} \path{GroupViewSystems.lean}, line 107.\par\noindent\path{KomorebiSpace.GroupViewSystem.twisted_coherent_iff}\par{\sffamily\footnotesize Source SHA-256 prefix: d2dd13e5f5cc2d78.}}\par
\par\smallskip\noindent{\small\textbf{F} \path{CocycleGaugeGroupoid.lean}, line 95.\par\noindent\path{KomorebiSpace.GroupViewSystem.gaugeSetoid}\par{\sffamily\footnotesize Source SHA-256 prefix: 93393322cd38df5e.}}\par
\smallskip\noindent{\color{KomorebiMuted}\rule{\linewidth}{.3pt}}\par\end{minipage}\par\medskip
\par\noindent\begin{minipage}{\linewidth}\subsection*{KS31 / Nonabelian extension obstruction}
A pointed based-gauge orbit set, not an assumed cohomology group or a general algorithm.
\par\smallskip\noindent{\small\textbf{F} \path{NonabelianExtension.lean}, line 96.\par\noindent\path{KomorebiSpace.GroupViewSystem.extensionOrbit_neutral_iff}\par{\sffamily\footnotesize Source SHA-256 prefix: 7834ff4193ad9a05.}}\par
\par\smallskip\noindent{\small\textbf{F} \path{NonabelianExtension.lean}, line 114.\par\noindent\path{KomorebiSpace.GroupViewSystem.extensionOrbit_independent}\par{\sffamily\footnotesize Source SHA-256 prefix: 7834ff4193ad9a05.}}\par
\smallskip\noindent{\color{KomorebiMuted}\rule{\linewidth}{.3pt}}\par\end{minipage}\par\medskip
\par\noindent\begin{minipage}{\linewidth}\subsection*{KS32 / Nonabelian extension torsor}
Parameters act by right multiplication on a chosen extension.
\par\smallskip\noindent{\small\textbf{F} \path{NonabelianExtensionFibers.lean}, line 51.\par\noindent\path{KomorebiSpace.GroupViewSystem.extensionParameters}\par{\sffamily\footnotesize Source SHA-256 prefix: 486c0bbabc34525c.}}\par
\par\smallskip\noindent{\small\textbf{F} \path{NonabelianExtensionFibers.lean}, line 107.\par\noindent\path{KomorebiSpace.GroupViewSystem.unique_extension_iff}\par{\sffamily\footnotesize Source SHA-256 prefix: 486c0bbabc34525c.}}\par
\smallskip\noindent{\color{KomorebiMuted}\rule{\linewidth}{.3pt}}\par\end{minipage}\par\medskip
\par\noindent\begin{minipage}{\linewidth}\subsection*{KS33 / S3 loop model}
Cocycle coherence is not the existence of a twisted global section.
\par\smallskip\noindent{\small\textbf{F} \path{CocycleExamples.lean}, line 205.\par\noindent\path{KomorebiSpace.CocycleExamples.s3_orbit_dichotomy}\par{\sffamily\footnotesize Source SHA-256 prefix: 6761e0443a1c1ad7.}}\par
\par\smallskip\noindent{\small\textbf{F} \path{CocycleExamples.lean}, line 169.\par\noindent\path{KomorebiSpace.CocycleExamples.actual_automorphisms_noncommute}\par{\sffamily\footnotesize Source SHA-256 prefix: 6761e0443a1c1ad7.}}\par
\par\smallskip\noindent{\small\textbf{F} \path{CocycleExamples.lean}, line 103.\par\noindent\path{KomorebiSpace.CocycleExamples.involution_trivial_iff}\par{\sffamily\footnotesize Source SHA-256 prefix: 6761e0443a1c1ad7.}}\par
\smallskip\noindent{\color{KomorebiMuted}\rule{\linewidth}{.3pt}}\par\end{minipage}\par\medskip
\par\noindent\begin{minipage}{\linewidth}\subsection*{KS34 / Effective groupoid descent}
Unframed gauge groupoids and composition-covered full patches. Uniqueness respects the specified identifications.
\par\smallskip\noindent{\small\textbf{F} \path{DescentEquivalence.lean}, line 104.\par\noindent\path{KomorebiSpace.GroupViewSystem.CocycleDescentData.descentEquivalence}\par{\sffamily\footnotesize Source SHA-256 prefix: c7ca73d617a578ed.}}\par
\par\smallskip\noindent{\small\textbf{F} \path{DescentComparison.lean}, line 79.\par\noindent\path{KomorebiSpace.GroupViewSystem.CocycleDescentData.Effectuation.exists_unique_comparison}\par{\sffamily\footnotesize Source SHA-256 prefix: afc5bf933284a648.}}\par
\smallskip\noindent{\color{KomorebiMuted}\rule{\linewidth}{.3pt}}\par\end{minipage}\par\medskip
\par\noindent\begin{minipage}{\linewidth}\subsection*{KS35 / Composition topology}
An actual Grothendieck topology, with the standard generated-sieve characterization.
\par\smallskip\noindent{\small\textbf{F} \path{CompositionSite.lean}, line 57.\par\noindent\path{KomorebiSpace.CompositionSite.topology}\par{\sffamily\footnotesize Source SHA-256 prefix: ac958295633fc50f.}}\par
\par\smallskip\noindent{\small\textbf{F} \path{CompositionSite.lean}, line 143.\par\noindent\path{KomorebiSpace.CompositionSite.ofArrows_covering_iff}\par{\sffamily\footnotesize Source SHA-256 prefix: ac958295633fc50f.}}\par
\smallskip\noindent{\color{KomorebiMuted}\rule{\linewidth}{.3pt}}\par\end{minipage}\par\medskip
\par\noindent\begin{minipage}{\linewidth}\subsection*{KS36 / Standard stackhood}
Target is mathlib Pseudofunctor.IsStack on the chosen composition site.
\par\smallskip\noindent{\small\textbf{F} \path{CompositionStack.lean}, line 178.\par\noindent\path{KomorebiSpace.GroupViewSystem.patch_isStack}\par{\sffamily\footnotesize Source SHA-256 prefix: 656bbbc1ddaba844.}}\par
\smallskip\noindent{\color{KomorebiMuted}\rule{\linewidth}{.3pt}}\par\end{minipage}\par\medskip
\par\noindent\begin{minipage}{\linewidth}\subsection*{KS37 / Equivalent view presentations}
Prescribed domains are isomorphism closed and coefficient systems are really pulled back.
\par\smallskip\noindent{\small\textbf{F} \path{SectionViewEquivalence.lean}, line 63.\par\noindent\path{KomorebiSpace.ViewSystem.viewEquivalence}\par{\sffamily\footnotesize Source SHA-256 prefix: 999bb5c576cefc7b.}}\par
\par\smallskip\noindent{\small\textbf{F} \path{RelativeSections.lean}, line 66.\par\noindent\path{KomorebiSpace.ViewSystem.relativeViewEquivalence}\par{\sffamily\footnotesize Source SHA-256 prefix: c28de5ab91b69828.}}\par
\par\smallskip\noindent{\small\textbf{F} \path{RelativeAmbiguity.lean}, line 43.\par\noindent\path{KomorebiSpace.GroupViewSystem.invisibleViewEquivalence}\par{\sffamily\footnotesize Source SHA-256 prefix: a24937a0816305f9.}}\par
\smallskip\noindent{\color{KomorebiMuted}\rule{\linewidth}{.3pt}}\par\end{minipage}\par\medskip
\par\noindent\begin{minipage}{\linewidth}\subsection*{KS38 / Identified-boundary realization groupoid}
Groupoid homotopy fiber construction. No new higher-category interface is asserted.
\par\smallskip\noindent{\small\textbf{F} \path{MarkedRealization.lean}, line 64.\par\noindent\path{KomorebiSpace.MarkedRealization.nonempty_iff}\par{\sffamily\footnotesize Source SHA-256 prefix: a3ad06525f28defd.}}\par
\par\smallskip\noindent{\small\textbf{F} \path{MarkedRealization.lean}, line 73.\par\noindent\path{KomorebiSpace.MarkedRealization.endomorphism_condition}\par{\sffamily\footnotesize Source SHA-256 prefix: a3ad06525f28defd.}}\par
\smallskip\noindent{\color{KomorebiMuted}\rule{\linewidth}{.3pt}}\par\end{minipage}\par\medskip
\par\noindent\begin{minipage}{\linewidth}\subsection*{KS39 / Marked transport across a square}
The square commutes through a specified natural isomorphism and the horizontal functors are equivalences.
\par\smallskip\noindent{\small\textbf{F} \path{MarkedTransport.lean}, line 117.\par\noindent\path{KomorebiSpace.MarkedRealization.transportEquivalence}\par{\sffamily\footnotesize Source SHA-256 prefix: 99fc0cf03cfdbdf9.}}\par
\par\smallskip\noindent{\small\textbf{F} \path{MarkedTransport.lean}, line 128.\par\noindent\path{KomorebiSpace.MarkedRealization.transportHomEquiv}\par{\sffamily\footnotesize Source SHA-256 prefix: 99fc0cf03cfdbdf9.}}\par
\smallskip\noindent{\color{KomorebiMuted}\rule{\linewidth}{.3pt}}\par\end{minipage}\par\medskip
\par\noindent\begin{minipage}{\linewidth}\subsection*{KS40 / Transported path telescope}
Known bounded invertible transports. Path gain is not replaced by length without contractivity.
\par\smallskip\noindent{\small\textbf{F} \path{AmplifiedTransport.lean}, line 151.\par\noindent\path{KomorebiSpace.AmplifiedTransport.Connection.walkSum_defect}\par{\sffamily\footnotesize Source SHA-256 prefix: da79596a2f437a07.}}\par
\par\smallskip\noindent{\small\textbf{F} \path{AmplifiedTransport.lean}, line 239.\par\noindent\path{KomorebiSpace.AmplifiedTransport.Connection.walk_error}\par{\sffamily\footnotesize Source SHA-256 prefix: da79596a2f437a07.}}\par
\smallskip\noindent{\color{KomorebiMuted}\rule{\linewidth}{.3pt}}\par\end{minipage}\par\medskip
\par\noindent\begin{minipage}{\linewidth}\subsection*{KS41 / Holonomy classification}
A bounded linear equivalence. It is an isometry only in the isometric subclass.
\par\smallskip\noindent{\small\textbf{F} \path{AmplifiedHolonomy.lean}, line 174.\par\noindent\path{KomorebiSpace.AmplifiedTransport.Rooted.parameterEquiv}\par{\sffamily\footnotesize Source SHA-256 prefix: f46a0ba9482bc51c.}}\par
\par\smallskip\noindent{\small\textbf{F} \path{AmplifiedHolonomy.lean}, line 104.\par\noindent\path{KomorebiSpace.AmplifiedTransport.Rooted.fixedSpace_route_independent}\par{\sffamily\footnotesize Source SHA-256 prefix: f46a0ba9482bc51c.}}\par
\smallskip\noindent{\color{KomorebiMuted}\rule{\linewidth}{.3pt}}\par\end{minipage}\par\medskip
\par\noindent\begin{minipage}{\linewidth}\subsection*{KS42 / Orthogonal exact repair}
Requires the declared Hilbert geometry, not closed residual range.
\par\smallskip\noindent{\small\textbf{M} \path{HilbertKernelProjection.lean}, line 245.\par\noindent\path{KomorebiSpace.HilbertKernelProjection.unconditional_exact_repair}\par{\sffamily\footnotesize Source SHA-256 prefix: 5a30a4de9c18ed34.}}\par
\par\smallskip\noindent{\small\textbf{M} \path{HilbertKernelProjection.lean}, line 120.\par\noindent\path{KomorebiSpace.HilbertKernelProjection.projection_best}\par{\sffamily\footnotesize Source SHA-256 prefix: 5a30a4de9c18ed34.}}\par
\par\smallskip\noindent{\small\textbf{M} \path{HilbertKernelProjection.lean}, line 134.\par\noindent\path{KomorebiSpace.HilbertKernelProjection.projection_unique}\par{\sffamily\footnotesize Source SHA-256 prefix: 5a30a4de9c18ed34.}}\par
\smallskip\noindent{\color{KomorebiMuted}\rule{\linewidth}{.3pt}}\par\end{minipage}\par\medskip
\par\noindent\begin{minipage}{\linewidth}\subsection*{KS43 / Closed range equivalence}
Banach target in the stated theorem. The explicit infinite diagonal illustration is written, not newly compiled.
\par\smallskip\noindent{\small\textbf{M} \path{HilbertKernelProjection.lean}, line 210.\par\noindent\path{KomorebiSpace.HilbertKernelProjection.closed_range_iff_bound}\par{\sffamily\footnotesize Source SHA-256 prefix: 5a30a4de9c18ed34.}}\par
\par\smallskip\noindent{\small\textbf{M} \path{HilbertKernelProjection.lean}, line 219.\par\noindent\path{KomorebiSpace.HilbertKernelProjection.closed_range_iff_near_exact}\par{\sffamily\footnotesize Source SHA-256 prefix: 5a30a4de9c18ed34.}}\par
\par\smallskip\noindent{\small\textbf{M} \path{HilbertKernelProjection.lean}, line 163.\par\noindent\path{KomorebiSpace.HilbertKernelProjection.gain_characterization}\par{\sffamily\footnotesize Source SHA-256 prefix: 5a30a4de9c18ed34.}}\par
\smallskip\noindent{\color{KomorebiMuted}\rule{\linewidth}{.3pt}}\par\end{minipage}\par\medskip
\par\noindent\begin{minipage}{\linewidth}\subsection*{KS44 / Amplified holonomy gain}
The finite-dimensional result is in the foundation archive. The infinite-dimensional necessity is in a separately recovered source, not claimed under the focused archive replay.
\par\smallskip\noindent{\small\textbf{F} \path{AmplifiedRepair.lean}, line 204.\par\noindent\path{KomorebiSpace.AmplifiedTransport.Rooted.orthogonalRepair_error}\par{\sffamily\footnotesize Source SHA-256 prefix: e4009d19b75ab9cf.}}\par
\par\smallskip\noindent{\small\textbf{S} \path{HilbertHolonomyRepair.lean}, line 194.\par\noindent\path{closed_range_iff_near_parallel}\par{\sffamily\footnotesize Source SHA-256 prefix: 5baf9b3a40b4fdd8.}}\par
\smallskip\noindent{\color{KomorebiMuted}\rule{\linewidth}{.3pt}}\par\end{minipage}\par\medskip
\par\noindent\begin{minipage}{\linewidth}\subsection*{KS45 / Flat transport amplification}
Both directed edge orientations use the original maximum norm.
\par\smallskip\noindent{\small\textbf{F} \path{AmplifiedExamples.lean}, line 139.\par\noindent\path{KomorebiSpace.AmplifiedTransport.Examples.exponential_amplification}\par{\sffamily\footnotesize Source SHA-256 prefix: e95779e1b481e79e.}}\par
\par\smallskip\noindent{\small\textbf{F} \path{AmplifiedExamples.lean}, line 148.\par\noindent\path{KomorebiSpace.AmplifiedTransport.Examples.length_only_bound_false}\par{\sffamily\footnotesize Source SHA-256 prefix: e95779e1b481e79e.}}\par
\smallskip\noindent{\color{KomorebiMuted}\rule{\linewidth}{.3pt}}\par\end{minipage}\par\medskip
\par\noindent\begin{minipage}{\linewidth}\subsection*{KS46 / Two-stage joint projection}
Second projection is inside the locally exact subspace.
\par\smallskip\noindent{\small\textbf{M} \path{QuantitativeDescent.lean}, line 61.\par\noindent\path{KomorebiSpace.QuantitativeDescent.repair_eq_projection}\par{\sffamily\footnotesize Source SHA-256 prefix: 0a35307bbdb5fdc3.}}\par
\par\smallskip\noindent{\small\textbf{M} \path{QuantitativeDescent.lean}, line 121.\par\noindent\path{KomorebiSpace.QuantitativeDescent.error_squared}\par{\sffamily\footnotesize Source SHA-256 prefix: 0a35307bbdb5fdc3.}}\par
\smallskip\noindent{\color{KomorebiMuted}\rule{\linewidth}{.3pt}}\par\end{minipage}\par\medskip
\par\noindent\begin{minipage}{\linewidth}\subsection*{KS47 / Joint closed range}
Closed local range is retained. The printed bound is a sufficient, not necessarily optimal, gain.
\par\smallskip\noindent{\small\textbf{M} \path{DescentClosedRange.lean}, line 93.\par\noindent\path{KomorebiSpace.QuantitativeDescent.joint_closed_iff_matching_closed}\par{\sffamily\footnotesize Source SHA-256 prefix: 394f9a40603290cd.}}\par
\par\smallskip\noindent{\small\textbf{M} \path{DescentClosedRange.lean}, line 70.\par\noindent\path{KomorebiSpace.QuantitativeDescent.stable_repair}\par{\sffamily\footnotesize Source SHA-256 prefix: 394f9a40603290cd.}}\par
\smallskip\noindent{\color{KomorebiMuted}\rule{\linewidth}{.3pt}}\par\end{minipage}\par\medskip
\par\noindent\begin{minipage}{\linewidth}\subsection*{KS48 / Two-sided projector motion}
Both independent residual certificates are supplied. No inclusion of kernels is assumed.
\par\smallskip\noindent{\small\textbf{M} \path{MovingKernelBounds.lean}, line 95.\par\noindent\path{KomorebiSpace.MovingKernel.projection_motion}\par{\sffamily\footnotesize Source SHA-256 prefix: 0edf4cbd9a1b1063.}}\par
\par\smallskip\noindent{\small\textbf{M} \path{MovingKernelBounds.lean}, line 164.\par\noindent\path{KomorebiSpace.MovingKernel.projection_lipschitz}\par{\sffamily\footnotesize Source SHA-256 prefix: 0edf4cbd9a1b1063.}}\par
\smallskip\noindent{\color{KomorebiMuted}\rule{\linewidth}{.3pt}}\par\end{minipage}\par\medskip
\par\noindent\begin{minipage}{\linewidth}\subsection*{KS49 / Derived perturbed certificate}
Banach target, supplied split surjectivity, and strict a times epsilon less than one.
\par\smallskip\noindent{\small\textbf{M} \path{SplitKernelPerturbation.lean}, line 92.\par\noindent\path{KomorebiSpace.SplitKernel.Certificate.perturbed_right_inverse}\par{\sffamily\footnotesize Source SHA-256 prefix: 1e60514ffacf125c.}}\par
\par\smallskip\noindent{\small\textbf{M} \path{SplitKernelPerturbation.lean}, line 98.\par\noindent\path{KomorebiSpace.SplitKernel.Certificate.perturbed_section_bound}\par{\sffamily\footnotesize Source SHA-256 prefix: 1e60514ffacf125c.}}\par
\par\smallskip\noindent{\small\textbf{M} \path{SplitKernelPerturbation.lean}, line 140.\par\noindent\path{KomorebiSpace.SplitKernel.Certificate.projection_motion}\par{\sffamily\footnotesize Source SHA-256 prefix: 1e60514ffacf125c.}}\par
\smallskip\noindent{\color{KomorebiMuted}\rule{\linewidth}{.3pt}}\par\end{minipage}\par\medskip
\par\noindent\begin{minipage}{\linewidth}\subsection*{KS50 / Transport between moving kernels}
Transport depends on the splitting and is not generally orthogonal projection.
\par\smallskip\noindent{\small\textbf{M} \path{SplitKernelPerturbation.lean}, line 216.\par\noindent\path{KomorebiSpace.SplitKernel.Certificate.kernelEquiv}\par{\sffamily\footnotesize Source SHA-256 prefix: 1e60514ffacf125c.}}\par
\par\smallskip\noindent{\small\textbf{M} \path{SplitKernelPerturbation.lean}, line 226.\par\noindent\path{KomorebiSpace.SplitKernel.Certificate.kernel_transport_error}\par{\sffamily\footnotesize Source SHA-256 prefix: 1e60514ffacf125c.}}\par
\smallskip\noindent{\color{KomorebiMuted}\rule{\linewidth}{.3pt}}\par\end{minipage}\par\medskip
\par\noindent\begin{minipage}{\linewidth}\subsection*{KS51 / Exact moving-line example}
The all-real projector formula is separate from the strict-radius certificate interface.
\par\smallskip\noindent{\small\textbf{M} \path{MovingKernelExamples.lean}, line 91.\par\noindent\path{KomorebiSpace.MovingKernelExamples.projection_formula}\par{\sffamily\footnotesize Source SHA-256 prefix: 8ce528290f97bfde.}}\par
\par\smallskip\noindent{\small\textbf{M} \path{MovingKernelExamples.lean}, line 142.\par\noindent\path{KomorebiSpace.MovingKernelExamples.motion_opNorm_exact}\par{\sffamily\footnotesize Source SHA-256 prefix: 8ce528290f97bfde.}}\par
\par\smallskip\noindent{\small\textbf{M} \path{MovingKernelExamples.lean}, line 199.\par\noindent\path{KomorebiSpace.MovingKernelExamples.kernel_transport_formula}\par{\sffamily\footnotesize Source SHA-256 prefix: 8ce528290f97bfde.}}\par
\smallskip\noindent{\color{KomorebiMuted}\rule{\linewidth}{.3pt}}\par\end{minipage}\par\medskip
\par\noindent\begin{minipage}{\linewidth}\subsection*{KS52 / Rank loss and radius endpoint}
The rank-changing diagonal pair is a written elementary example. Source correspondence covers the general loss estimate and the radius boundary.
\par\smallskip\noindent{\small\textbf{M} \path{MovingKernelBounds.lean}, line 148.\par\noindent\path{KomorebiSpace.MovingKernel.lost_exact_unit_forces_gain}\par{\sffamily\footnotesize Source SHA-256 prefix: 0edf4cbd9a1b1063.}}\par
\par\smallskip\noindent{\small\textbf{M} \path{MovingKernelExamples.lean}, line 227.\par\noindent\path{KomorebiSpace.MovingKernelExamples.closed_radius_failure}\par{\sffamily\footnotesize Source SHA-256 prefix: 8ce528290f97bfde.}}\par
\smallskip\noindent{\color{KomorebiMuted}\rule{\linewidth}{.3pt}}\par\end{minipage}\par\medskip
\par\noindent\begin{minipage}{\linewidth}\subsection*{KS53 / Motion of quantitative descent}
Uses the original two-stage repair through its proved identification with the joint projector.
\par\smallskip\noindent{\small\textbf{M} \path{MovingDescent.lean}, line 52.\par\noindent\path{KomorebiSpace.MovingDescent.repair_motion}\par{\sffamily\footnotesize Source SHA-256 prefix: 6820862031ad7094.}}\par
\smallskip\noindent{\color{KomorebiMuted}\rule{\linewidth}{.3pt}}\par\end{minipage}\par\medskip
\par\noindent\begin{minipage}{\linewidth}\subsection*{KS54 / Finite prefixes without a point}
Additional written counterexample with a complete elementary proof. No new Lean declaration or universal computational proof is claimed.
\par\smallskip\noindent\textbf{W.} The proof is written in the monograph. No new Lean declaration or compiler verification is claimed.
\smallskip\noindent{\color{KomorebiMuted}\rule{\linewidth}{.3pt}}\par\end{minipage}\par\medskip
